\documentclass[quick,pro]{debate}
\motion{「勇敢的人先享受世界」是一个陷阱 / 不是一个陷阱}
\stance{正方}{是一个陷阱}
\begin{document}
\debatetitle
\begin{thesisbox}它把「有没有兜底」这张门票，标成了「够不够勇敢」这枚勋章，于是最付不起代价的人被叫去付账。\end{thesisbox}

\section{定义与判准（标准表述）}
\begin{fields}
\field{陷阱}{它呈现出来的样子和照做之后的真实后果系统性地不一致，而且正是这层外观把人引了进去。\sub{关键}{陷阱不需要有人挖，中等收入陷阱就没有设局者。}}
\field{勇敢的人}{在后果不确定时仍然选择行动的人。不等于鲁莽，也不等于有钱。}
\field{先享受世界}{比不敢行动的人更早得到世界给的好处。}
\field{判准}{三件事：① 外观与真实因果系统性不一致；② 照做的人受损是常态；③ 受损正是被它遮住的东西造成的。\sub{我方要证明}{三件都成立，举证责任在我方。}\sub{对方要证明}{打掉任意一件即可。}}
\field{顶替线（全场安全带）}{省略条件只是不严谨；拿另一个变量顶替被省略的那个，才是陷阱。「早睡早起身体好」没告诉你睡不好是因为你不自律，这句话告诉你了。}
\field{对方提偏向定义或判准时的 30 秒口径}{要「必须有设局者」：中等收入陷阱是谁设的局？要「必须有受损统计」：这句话才流行三年，谁都没有，那就比机制。}
\end{fields}

\section{我方论点}
\begin{tabularx}{\linewidth}{@{}L{4.6em}Y Y L{7em}@{}}\toprule
\thead{标签}&\thead{一句话}&\thead{核心数据 / 学理 / 案例}&\thead{出处}\\\midrule
\textbf{门票换勋章}&它把「有没有兜底」这张门票，标成了「够不够勇敢」这枚勋章&数据：23.8\% 不 GAP 因经济压力，另 18.5\% 怕简历断档；学理：基本归因错误；案例：「勇敢」从冲动变成物质实力&智联招聘 2025 春季跳槽调研，37980 份样本；Ross 1977；澎湃 2023-10-17\\
\textbf{删过的账本}&它给的是一份删掉了失败样本、也删掉了反面证据的账本&数据：长期回顾 51\% 困扰于做过的事，49\% 困扰于没做的事；学理：幸存者偏差；案例：不含在校生青年失业率 14.9\%&Richardson \& Gilovich 2023, n=2600；Wald 1943；国家统计局 2026 年 6 月\\
\bottomrule\end{tabularx}

\section{对方论点 → 一句话反驳}
\begin{tabularx}{\linewidth}{@{}Y Y Y@{}}\toprule
\thead{对方会说}&\thead{我方一句话回}&\thead{对方追问时}\\\midrule
它给的是顺序，不是支票&顺序也是承诺，它承诺的是「先」；原句是「先享受」，不是「先开始」&问「盖住的是哪几个字」→ 兜底能力\\
摔的是没留退路，不是有勇气&那你方守的是「准备好的人先享受世界」&问「词典义不算吗」→ 算，勇气针对已知风险，而它不提供风险信息\\
陷阱必须有设局者&中等收入陷阱是谁设的局？&问「那什么不是陷阱」→ 说了条件的经验分享不是，只省略没顶替的格言不是\\
1994 年研究说人更后悔没做的事&请把 2023 年那份 2600 人的复制念完，51\% 对 49\%&说「五五开不算陷阱」→ 五五开的是后悔，不是代价\\
\bottomrule\end{tabularx}

\section{必问与必答}
\begin{points}{必问（对方不答就点名、重复、不换题）}
\item 这九个字里，哪个字提到了存款？（主问题）
\item 「享受」两个字，对方删到哪里去了？
\item 陷阱需不需要有设局者？请给一个字的回答。
\end{points}
\begin{points}{必答（15 秒正面回答）}
\item 「照它做的人普遍变差，数据呢？」→ 没有直接数据，这话才流行三年；我方打代价错配，对方也没有反面数据。
\item 「那什么不是陷阱？」→ 说了条件的经验分享不是，只省略、没顶替的格言不是。
\item 「你方是不是在劝人别勇敢？」→ 不是。先看清门票价格再决定买不买。
\end{points}

\section{自由辩战场概要}
\begin{tabularx}{\linewidth}{@{}L{4.2em}Y Y Y@{}}\toprule
\thead{战场}&\thead{开场问题}&\thead{目标}&\thead{收束语}\\\midrule
门票（前 90 秒）&哪个字提到了存款？&决定能不能走的是条件，说出口的却是勇敢&门票在，标签换了\\
账本（中 90 秒）&2023 年那一半，对方念不念？&说服力建立在一个已被推翻的常识上&做和不做一样，凭什么只劝人做\\
\bottomrule\end{tabularx}
时间：前 90 秒战场一并抛完必问，中 90 秒战场二，后 60 秒只重复主问题、不开新战场。站位：一辩守定义判准，二辩接驳论，三辩接盘问，四辩收束。

\textbf{各环节}：1 正一立论 180 秒，二十秒内把三件事说完；3 正四质询反一 120 秒双边，拿到「陷阱不需要设局者」；4 正二驳论 120 秒，留临场位；6 正三盘问 90 秒，先问反二再问反四同一问题；10 正四结辩 210 秒，留 30–40 秒真回应反四。被质询与被盘问时每答不超过八秒，不否认常识。

\section{如果……就……}
\begin{contingency}
\ifthen{对方定义要「必须有设局者」}{中等收入陷阱是谁设的局？请不要用定义取消比赛。}
\ifthen{对方把我方打成「什么都叫陷阱」}{三秒内背出顶替线：省略不算，顶替才算。}
\ifthen{「门票换勋章」被打穿}{收缩到「删过的账本」，用判准第一件解释为什么仍然占优。}
\ifthen{对方回避必问或拿出没预料到的数据}{回避就点名、重复、不换题；新数据先问出处与自变量。}
\end{contingency}

\end{document}
